\section{Certified vector intertwiners and fast unknown-base recovery}
\label{sec:vector-algorithm}
The symmetric full-kernel algorithm also handles nonsymmetric quotient
matrices. The reduction computes a doubled kernel, then projects to its
right component before enumeration. Its expected-work application needs
one bounded moment for that right kernel in the actual graph law.

\subsection{A complete right kernel on every input}
Fix an odd prime $p$. Let $A=A^T$ be an integer adjacency matrix of a
simple graph of maximum degree at most $D$, and let $R$ be a fixed
$r\times r$ integer matrix with bounded entries. Stack matrix columns
in the convention $\operatorname{vec}(X)=(X_{:,1},\ldots,X_{:,r})$.
Over $\Fp$, put
\begin{equation}\label{eq:vector-operator}
 \mathcal K_R=I_r\otimes A-R^T\otimes I_n,\qquad
 m=rn,\quad w=\operatorname{nnz}(\mathcal K_R),\quad
 h=\nullity_{\Fp}\mathcal K_R.
\end{equation}
Then $\mathcal K_R\operatorname{vec}(X)=\operatorname{vec}(AX-XR)$;
its transpose acts as $\operatorname{vec}(Y)\mapsto
\operatorname{vec}(AY-YR^T)$. Each product costs
$O(r\operatorname{nnz}(A)+n\operatorname{nnz}(R))$ field operations.
At fixed $r,D$ and bounded quotient entries, $m=O(n)$ and $w=O(n)$.
The structured products or an explicit sparse representation suffice.

\begin{theorem}[Certified right kernels]\label{vec:thm:kernel}
For every fixed input in \eqref{eq:vector-operator}, a Las Vegas
algorithm returns the complete right-kernel basis
$C\in\Fp^{m\times h}$, normalized on pivot rows $J$ by $C_{J,:}=I_h$.
Its conditional expected bit work is
\[
 \widetilde O_p\bigl(m(w+m)(h+1)\bigr),
\]
and its working storage is
$\widetilde O_p(m(h+1)+w)$ bits, apart from a retry counter.
A certificate of $O_p(mh+m\log m)$ bits is deterministically
checkable in $\widetilde O_p(m(w+m))$ bit work.
Every accepted answer is correct, including on inputs with large
nullity. No symmetry assumption on $R$ is imposed.
\end{theorem}
\begin{proof}
Apply Theorem~\ref{fk:thm:kernel} to the symmetric matrix
\begin{equation}\label{eq:vector-doubling}
 \mathcal S_R=
 \begin{pmatrix}0&\mathcal K_R\\\mathcal K_R^T&0\end{pmatrix}
 \quad\text{of order }2m.
\end{equation}
For $(u,v)$ it acts as $(\mathcal K_Rv,\mathcal K_R^Tu)$, so
\[
 \ker\mathcal S_R=\ker\mathcal K_R^T\oplus\ker\mathcal K_R,
 \qquad \rank\mathcal S_R=2\rank\mathcal K_R,
 \qquad \nullity\mathcal S_R=2h.
\]
These are field identities; no Gram-matrix or positivity argument is
used. One product uses one $\mathcal K_R$ and one $\mathcal K_R^T$
product. The per-input theorem, with order $2m$, sparsity $2w$, and
nullity $2h$, has the displayed work and storage bounds. Its verified
extension-field construction, preconditioning retries and base-field
descent are all included.

The bottom $m$ rows of a complete $\mathcal S_R$-kernel basis span
$\ker\mathcal K_R$: every right nullvector $v$ occurs as $(0,v)$.
Select $h$ independent projected columns by exact elimination, choose
$h$ invertible pivot rows, and normalize there. Arbitrary first columns
need not be independent. Selection and normalization cost
$O(mh^2+h^3)\le O(m^2(h+1))$, within the theorem's work bound.
They use $O(mh)$ storage. The symmetric primitive has already returned
a base-field basis. Equivalently, a projected extension-field basis
can be normalized first: Frobenius preserves the base-defined right
kernel and its uniquely pivot-normalized basis, so that basis lies in
$\Fp$. Extract its base coefficients and check $\mathcal K_RC=0$.
The cases $h=0$ and $h=m$ require no exception to these arguments.

For a compact final certificate, retain the symmetric primitive's
verified field definition, nonzero diagonal preconditioner $\mathcal D$,
and scalar probes $u,v$, together with $C,J$. Its checker regenerates
the $4m$ terms
\[
 u^T(\mathcal D\mathcal S_R)^{j+1}v,
 \qquad 0\le j<4m,
\]
and recomputes their true shortest Berlekamp--Massey length $\ell$.
Require $0\le h\le m$, the prescribed dimensions, distinct pivot
rows $J$, and canonical base-field entries of $C$. Require
$\ell=2(m-h)$ and check $\mathcal K_RC=0$ and $C_{J,:}=I_h$.
The unconditional shifted-sequence bound of Lemma~\ref{fk:lem:lower}
gives $\ell\le\rank\mathcal S_R=2\rank\mathcal K_R$ on every draw.
The verified independent columns give $\rank\mathcal K_R\le m-h$.
Thus acceptance forces
\[
 m-h=\ell/2\le\rank\mathcal K_R\le m-h,
\]
certifying completeness with zero verification error. This is an
integer dimension comparison, not division by two in the field.
The verifier checks field sizes and encodings before the primitive's
irreducibility and sequence checks. Its $O(m)$ sparse products,
quadratic recurrence computation, and residual/pivot checks cost
$\widetilde O_p(m(w+m))$. Retaining the right basis and probes gives
the stated certificate size; the unused left basis is unnecessary.
\end{proof}

\subsection{Enumeration under a single actual-law moment}
\begin{theorem}[Moment-to-vector-list theorem]\label{vec:thm:moment}
Fix $K,D$ and a finite family $\mathcal R$ of nonnegative integer
quotient matrices of orders at most $K$, with row sums at most $D$.
Suppose a law on simple graphs of maximum degree at most $D$ satisfies,
for one fixed odd prime $p$ and every $R\in\mathcal R$,
\begin{equation}\label{eq:vector-moment}
 \sup_n\E p^{h_R(A)}\le C_R<\infty,\qquad
 h_R(A)=\nullity_{\Fp}(I_r\otimes A-R^T\otimes I_n).
\end{equation}
There is an always-correct Las Vegas algorithm which lists all
nonempty labeled equitable partitions with quotient in $\mathcal R$.
Its unconditional expected bit work is $\widetilde O(n^2)$, its
expected working storage is $\widetilde O(n)$ bits, and its expected
output size is $O(n)$ bits. Constants may depend on
$p,K,D,\mathcal R$ and the moment bounds. The algorithm is complete
on every input; the moment assumption controls its expected cost.
\end{theorem}
\begin{proof}
For each $R$, compute the certified right basis $C$ from
Theorem~\ref{vec:thm:kernel}. Enumerate only its $p^h$ base-field
coefficient vectors. Maintain the current vector using a base-$p$
odometer or a Gray order. Across all states there are $O_p(p^h)$
digit changes, each adding a scalar multiple of one column of $C$
in $O(m)$ field operations. Total generation costs
$\widetilde O_p(mp^h)$ bit work, with no extra factor $h$.
The doubled kernel is never enumerated.

Reshape each vector into $X\in\Fp^{n\times r}$. Require that every
row has exactly one entry equal to one and all remaining entries
zero, and that every column is nonzero. Lift these rows to integer
zero-one rows and check $AX=XR$ over the integers. Check any class
sizes actually prescribed by the task. A blind unknown-frame search
has no such size promise and imposes only nonemptiness. At fixed
$K,D$, these scans and sparse integer neighbor-count checks cost
$\widetilde O(n)$ per candidate. Every true integer partition reduces
to a field-kernel vector, and the integer check rejects modular
artifacts. Hence the resulting list is exact on every input.

For a fixed input and target the conditional expected work is
\begin{equation}\label{eq:vector-per-input}
 \widetilde O_p\bigl(m(w+m)(h+1)+mp^h\bigr).
\end{equation}
Because $h+1\le p^h$ and $m,w=O_{K,D}(n)$, assumption
\eqref{eq:vector-moment} bounds its graph expectation by
$\widetilde O(n^2)$. This averages the valid per-input Las Vegas
bound, without excluding rare large kernels. Neither
$\E p^{2h}$ nor $\E[h p^h]$ is used.
The same bound gives expected $\widetilde O(n)$ working storage.
Stream the output: at most $p^h$ candidates pass, and a partition
uses $O_K(n)$ bits. Its expected output size is therefore $O(n)$.
Finite summation over $\mathcal R$ preserves all these bounds;
independence of different runs is unnecessary. Reuse their workspace.
\end{proof}

For unknown quotients, use every order $1\le r\le K$ and all
nonnegative integer matrices with row sums at most $D$. There are
at most $\sum_{r=1}^K(D+1)^{r^2}$ candidates. Retain a verified
partition with the maximum number of classes, using only $O(n)$
additional bits. On the rigidity event of
Theorem~\ref{thm:unknown-frame}, every candidate coarsens the plant,
so this maximum is the planted partition, up to class names.
Completeness and zero algorithmic error hold outside that event too;
statistical recovery there is a separate question.

Nonsymmetric quotient candidates must remain in this blind search.
For any accepted one-hot $X$ with class sizes $N_i>0$, symmetry of
$A$ gives the integer balance identities
\[
 X^TAX=\diag(N_i)R=(\diag(N_i)R)^T,
 \qquad N_iR_{ij}=N_jR_{ji}.
\]
Equal positive sizes, when prescribed, force $R=R^T$. Equal sizes
of the true plant do not force its permitted coarsenings to have
equal sizes, so they do not justify discarding nonsymmetric candidates.
There is no division by a class size, including when $p\mid N_i$.
For a fixed $R$, only its allowed label automorphisms preserve that
target. To list unlabeled partitions in the unknown-$R$ search,
canonicalize $(X,R)$ simultaneously under the at most $K!$ label
permutations and output only the canonical representative. This
uses a constant number of scans and no growing deduplication table.

\subsection{Unknown certified frames in the actual simple law}
\begin{corollary}[Nearly quadratic unknown-frame recovery]
\label{vec:cor:frame}
Under the omission certificate of
Theorem~\ref{thm:unknown-frame}, the algorithm given only $A,K,D$
lists all equitable partitions with at most $K$ classes in
unconditional expected $\widetilde O_{K,D}(n^2)$ bit work, expected
$\widetilde O_{K,D}(n)$ working bits, and expected $O_{K,D}(n)$
output bits. Selecting a largest-class partition recovers the plant
and its unknown frame with failure $O_{K,D}(n^{-(d_{\min}-2)})$.
Varying total degrees, zero omission weights, internal degrees zero or
one, unequal planted sizes, and all
admissible sizes at which the simple factors exist are included.
\end{corollary}
\begin{proof}
Connectedness and integer balance fix the planted size ratios.
They can be written $N_i=c_iN$ for a primitive positive integer vector
$(c_i)$; this is only a parametrization of all admissible sizes.
Theorem~\ref{thm:frame-bounded-moment} gives
\eqref{eq:vector-moment} in exactly the independent uniform-simple
factor law, for $p=3$ and every reduction of every candidate $R$.
There are finitely many permitted planted frames, size-ratio vectors,
and candidate targets at fixed $K,D$, so a common moment constant
exists. Apply Theorem~\ref{vec:thm:moment} and the structural failure
bound of Theorem~\ref{thm:unknown-frame}.
\end{proof}

There is also a deterministic polynomial truncation preserving the
structural failure rate. Over $\F_2$, use dense elimination for each
candidate and cap its kernel size at $n^{d_{\min}-1}$, reading the minimum degree from the input. The bounded binary moment and a finite union give
interruption probability $O_{K,D}(n^{-(d_{\min}-1)})$. Gray enumeration and
integer checks then use bounded
$\widetilde O_{K,D}(n^3+n^{d_{\min}})$ bit work and $O_K(n^2)$ space.
Adding the structural failure preserves $O(n^{-(d_{\min}-2)})$.
The expected nearly quadratic Las Vegas algorithm has no such cap.

\subsection{The random-lift specialization}
\begin{corollary}[Nearly quadratic unknown-base recovery]
\label{vec:cor:lift}
Let $H$ be a fixed connected simple $d$-regular graph, $d\ge3$,
with $k\le K$ vertices, and let $A_N$ be its random $N$-lift,
with independent uniform matchings above all base edges. Given only
the lifted graph and fixed bounds $K,D\ge d$, the algorithm recovers
the base fibers and $H$ up to isomorphism with probability $1-o(1)$.
Its unconditional expected bit work is
$\widetilde O_{H,K,D}(n^2)$, where $n=kN$, with expected
$\widetilde O_{H,K,D}(n)$ working bits. Listing all equitable
partitions with at most $K$ classes has expected $O_{H,K,D}(n)$
output bits. All positive lift sizes are covered, including bipartite
bases; no real-corank or eigengap promise is needed.
\end{corollary}
\begin{proof}
Choose one fixed odd prime, for example $p=3$. The bounded lift
moment theorem, Theorem~\ref{thm:lift-bounded-moment}, states that
for this actual unconditioned lift law, every fixed $r$ and target
$\overline R\in\Fp^{r\times r}$ satisfy
\[
 \sup_N\E\bigl|\{Y\in\Fp^{n\times r}:A_NY=Y\overline R\}\bigr|
 <\infty.
\]
Apply it to the reductions of the finite integer quotient family
above. Its solution count is exactly $p^{h_R(A_N)}$, so
Theorem~\ref{vec:thm:moment} gives the work, space and output bounds.
A lift of a simple base is automatically simple; there is no
conditioning or change of graph law. Both $H$ and $R$ are analysis
parameters in this moment invocation, not information supplied to
the algorithm. Corollary~\ref{cor:random-lift} supplies the rigidity
event, and maximum-class selection recovers the plant on that event.
\end{proof}

The omission certificate is sufficient; minimum degree three alone is
not substituted for it. The nearly quadratic work is an expected bound,
not a worst-case quadratic guarantee.
